% ============================================================
% Model: Peierls 相变（一维二聚化）
% ============================================================

\section{Peierls 相变（一维二聚化）}
\label{sec:peierls}

\begin{tipbox}[关键词标签]
Peierls 不稳定性 · 一维二聚化 · 能带折叠 · 金属--绝缘体转变 · 派尔斯相变 · 电子--晶格耦合
\end{tipbox}

% --------------------------------------------------------
% 1. 模型定义框
% --------------------------------------------------------
\begin{modelbox}[模型：一维链的 Peierls 不稳定性]
\textbf{物理情境}：一维半满金属链在低温下不稳定，原子会发生位移形成周期加倍的二聚化结构（长键--短键交替），打开能隙，使系统从导体变为绝缘体。这是电子--晶格耦合导致的一维金属的普遍不稳定性。

\textbf{二聚化势场}：原子位移导致势场周期性从 $a$ 变为 $2a$：
\[
V(x)=V_0\cos\frac{\pi x}{a}=V_0\cos(q_0 x),\qquad q_0=\frac{\pi}{a}=2k_F
\]

\textbf{关键参数}：
\begin{itemize}[nosep]
    \item $a$：原晶格常数（单原子链周期）
    \item $2a$：二聚化后超晶格周期
    \item $k_F=\pi/(2a)$：半满一维链的费米波矢
    \item $q_0=2k_F=\pi/a$：二聚化波矢（恰好连接费米面上 $+k_F$ 与 $-k_F$）
    \item $\Delta$：二聚化打开的能隙，$\Delta\propto u$（原子位移量）
    \item $\lambda$：电子--晶格耦合常数
\end{itemize}

\textbf{物理图像}：
\begin{itemize}[nosep]
    \item 未二聚化：周期 $a$，第一布里渊区 $[-\pi/a,\pi/a]$，能带半满，金属
    \item 二聚化：周期 $2a$，布里渊区折叠为 $[-\pi/2a,\pi/2a]$，在 $k_F$ 处打开能隙，绝缘体
\end{itemize}
\end{modelbox}

% --------------------------------------------------------
% 2. 关键公式框
% --------------------------------------------------------
\begin{formulabox}[核心公式]
\begin{enumerate}[label=\textbf{(\arabic*)}]
    \item \textbf{一维半满链费米波矢}：
    \[
    k_F=\frac{\pi}{2a}
    \]
    \texteq{每个原子一个电子，$k$ 空间填满 $[-k_F,k_F]$，总长 $2k_F=\pi/a$}

    \item \textbf{二聚化后费米面落在能隙}：
    \[
    E_g=2\Delta=2|V_0|\propto u
    \]
    \texteq{$u$ 为原子位移振幅，能隙正比于位移}

    \item \textbf{电子能量增益（每个原子）}：
    \[
    \Delta E_{\text{el}}\approx-\frac{\Delta^2}{4\varepsilon_F}N(0)\cdot 2\Delta\sim-\Delta^2\ln\frac{\varepsilon_F}{\Delta}
    \]
    \texteq{对数发散的一维态密度使能量增益总是超过晶格弹性能}

    \item \textbf{晶格弹性能（每个原子）}：
    \[
    \Delta E_{\text{latt}}=\frac{1}{2}Ku^2
    \]
    \texteq{$K$ 为晶格恢复力常数，二次正比于位移}

    \item \textbf{Peierls 定理}：在一维电子--晶格耦合系统中，\textbf{任何} $T=0$ 的金属态都不稳定，总会发生二聚化打开能隙。
    \texteq{二维/三维系统中，电子能量增益 $\propto\Delta^2$ 不足以补偿弹性能，金属可稳定}
\end{enumerate}
\end{formulabox}

% --------------------------------------------------------
% 3. 训练点框
% --------------------------------------------------------
\begin{drillbox}[训练点与考法变体]
\trainingpoint{1} \textbf{能带折叠图像}：画出二聚化前后的 $E(k)$ 图，标出 $k_F$、布里渊区边界、能隙位置，说明为何费米面恰好落在能隙中。

\trainingpoint{2} \textbf{导电性变化分析}：二聚化前为半满金属，二聚化后变为绝缘体/半导体。用能带论解释电阻率的变化。

\trainingpoint{3} \textbf{与交替力常数链的对比}：Peierls 二聚化是``自发''的（电子能量驱动），而交替力常数链是``给定''的（外生结构）。两者在能谱上的异同。

\trainingpoint{4} \textbf{聚乙炔中的应用}：聚乙炔 $(\text{CH})_n$ 的 Peierls 二聚化导致键长交替（单键/双键），产生能隙，是导电聚合物的基础。

\trainingpoint{5} \textbf{高温恢复}：温度升高时，热涨落破坏二聚化序，体系恢复金属性。估计相变温度 $T_P\sim\Delta/k_B$。
\end{drillbox}

% --------------------------------------------------------
% 4. 解题技巧框
% --------------------------------------------------------
\begin{tipbox}[快速识别与解题技巧]
\begin{itemize}
    \item \examnote{识别标志}：题干出现``一维半满''''二聚化''''Peierls''''金属--绝缘体转变''''键长交替''，立即定位本模型。
    \item \examnote{布里渊区折叠速记}：
    \begin{center}
    \begin{tabular}{lll}
    \toprule
    \textbf{参数} & \textbf{二聚化前} & \textbf{二聚化后} \\
    \midrule
    周期 & $a$ & $2a$ \\
    布里渊区 & $[-\pi/a,\pi/a]$ & $[-\pi/2a,\pi/2a]$ \\
    倒格矢 & $G_n=2\pi n/a$ & $G_n'=\pi n/a$ \\
    $k_F$ & $\pi/2a$（恰为新边界） & 落在能隙中 \\
    \bottomrule
    \end{tabular}
    \end{center}
    \item \examnote{关键巧合}：一维半满链恰好 $k_F=\pi/2a$，二聚化后新的布里渊区边界就是 $k_F$。这是 Peierls 不稳定性在\textbf{任意}填充的一维链中都成立的根源（$q_0=2k_F$ 总是把费米面映射到自身）。
    \item \examnote{一维的特殊性}：一维费米面只有两点 $\pm k_F$，$2k_F$ 散射恰好连接这两点，产生发散响应。二维/三维的费米面是曲线/曲面，$2k_F$ 散射不集中在单一点，无发散。
    \item \examnote{能量竞争}：电子能量增益 $\sim\Delta^2\ln(\varepsilon_F/\Delta)$（一维对数发散），晶格弹性能 $\sim\Delta^2$。对数项使电子总能赢。
\end{itemize}
\end{tipbox}

% --------------------------------------------------------
% 5. 易错点框
% --------------------------------------------------------
\begin{errorbox}{常见失误与陷阱}
\begin{itemize}
    \item \textbf{费米波矢计算}：一维链，每个原子一个电子，自旋简并度2，$N=2\cdot(2k_F)/(2\pi/a)=2ak_F/\pi$（每原子）。故 $k_F=\pi/(2a)$。不要误用 $k_F=\pi/a$（那是每个原子\textbf{两个}电子的情况）。
    \item \textbf{能隙位置}：二聚化打开的能隙在\textbf{新的}布里渊区边界 $k=\pm\pi/2a$（即原 $k_F$ 处），不是在原布里渊区边界 $k=\pm\pi/a$。
    \item \textbf{二维/三维的稳定性}：Peierls 定理\textbf{只适用于一维}。二维/三维的金属可以稳定存在（如 Cu、Na），因为费米面的嵌套不完美，电子能量增益不足以补偿弹性能。
    \item \textbf{二聚化与双原子链混淆}：双原子链是``两种不同原子''或``两种不同力常数''；Peierls 二聚化是``相同原子因电子作用自发位移''。前者是结构给定，后者是电子驱动的自发对称性破缺。
    \item \textbf{SSH 模型}：Su--Schrieffer--Heeger 模型是 Peierls 二聚化的晶格模型，包含电子跃迁与晶格位移的耦合。不要与紧束缚模型混淆。
\end{itemize}
\end{errorbox}

% --------------------------------------------------------
% 6. 物理图像与关联模型
% --------------------------------------------------------
\begin{insightbox}[物理图像与模型关联]
\textbf{物理图像}：

想象一条由 identical 椅子排成的直线剧场：
\begin{itemize}[nosep]
    \item \textbf{正常态}：椅子等间距排列，观众（电子）可以自由走动，是``金属''。
    \item \textbf{Peierls 二聚化}：观众发现如果椅子两两靠近（形成``双人座''），整体能量更低——因为每对观众可以面对面交流（Cooper 对式的关联）。剧场于是变成了``双人座--宽过道''交替的结构。
    \item \textbf{能隙}：宽过道成了观众走动的``禁区''（能隙），只有能量足够高的观众才能跨过。
\end{itemize}

\textbf{与相似模型的对比}：
\begin{center}
\begin{tabular}{llll}
\toprule
\textbf{对比维度} & \textbf{Peierls 二聚化} & \textbf{CDW} & \textbf{超导体} \\
\midrule
维度 & 1D 为主 & 2D/3D（部分 nesting） & 任意 \\
序参量 & 晶格位移 $u$ & 电荷密度调制 $\rho(\mathbf{r})$ & 复数 $\psi=|\psi|e^{i\phi}$ \\
能隙 & 单粒子能隙 & 单粒子能隙 & 准粒子能隙 \\
驱动力 & 电子--晶格耦合 & 电子--电子相互作用 & 电子--声子耦合 \\
\bottomrule
\end{tabular}
\end{center}

\textbf{推广方向}：
\begin{itemize}[nosep]
    \item 聚乙炔中的孤子：二聚化域壁，携带分数电荷，是拓扑激发
    \item 电荷密度波（CDW）：2D/3D 中费米面 nesting 导致的类似现象（如 NbSe$_2$）
    \item 自旋派尔斯：自旋--晶格耦合导致的一维自旋链二聚化
\end{itemize}
\end{insightbox}

% --------------------------------------------------------
% 7. 推导附录
% --------------------------------------------------------
\begin{derivationbox}[推导细节（自我检验用）]
\textbf{Step 1：一维半满链的费米波矢}

$k$ 空间态密度（含自旋）：$L/\pi$。电子总数 $N$（$N$ 个原子，每个贡献1个电子）：
\[
N=\frac{L}{\pi}\cdot 2k_F=\frac{Na}{\pi}\cdot 2k_F\quad\Rightarrow\quad k_F=\frac{\pi}{2a}.
\]

\textbf{Step 2：二聚化后的能带}

原子位移 $u_n=(-1)^n u$，势场调制 $V(x)=V_0\cos(\pi x/a)$。

在 $k=\pm\pi/2a$ 附近，自由电子态 $|k\rangle$ 与 $|k-G\rangle$（$G=\pi/a$）简并。

$2\times2$ 哈密顿量：
\[
H=\begin{pmatrix}\varepsilon_k & \Delta\\ \Delta & \varepsilon_{k-G}\end{pmatrix},
\qquad \Delta=\frac{V_0}{2}.
\]

在 $k=\pi/2a$ 处，$\varepsilon_k=\varepsilon_{k-G}=\hbar^2\pi^2/(8ma^2)\equiv\varepsilon_F$。

本征值：$E=\varepsilon_F\pm\Delta$。能隙 $E_g=2\Delta$。

\textbf{Step 3：电子能量增益}

$k=\pi/2a$ 以下所有态能量降低。近似计算（线性化色散）：
\[
\Delta E_{\text{el}}=-\frac{L}{\pi}\int_{-k_F}^{k_F}\bigl[\sqrt{\varepsilon^2+\Delta^2}-|\varepsilon|\bigr]\,dk
\approx-\frac{L}{\pi v_F}\Delta^2\ln\frac{\varepsilon_F}{\Delta}.
\]

一维态密度的对数发散使能量增益总是大于弹性能 $\frac{1}{2}Ku^2\propto\Delta^2$。

\textbf{Step 4：为什么不是任意 $q$}

只有 $q_0=2k_F$ 的晶格畸变才能连接费米面上 nesting 的态（$+k_F$ 与 $-k_F$），使电子能量最大化降低。其他 $q$ 只打开远离费米面的能隙，电子能量增益很小。
\end{derivationbox}

\vspace{1em}
